\documentclass[12pt]{article}

\usepackage[utf8]{inputenc}
\usepackage[spanish]{babel}
\usepackage[shortlabels]{enumitem}

\usepackage{mathtools}
\usepackage{amssymb}
\usepackage{amsthm}

\usepackage{geometry}
\geometry{margin=2.5cm}

\usepackage{hyperref}

\begin{document}
\section*{Ejercicio 1}

Indica el orden y el grado de las siguientes ecuaciones diferenciales:

\begin{enumerate}[a)]
\item $(y'')^2 \sen(x) + 3x\,y\,y' + y^2 e^x = 0$
\item $y'' + 5x^2 (y')^3 + 6y = 0$
\item $x^2 \sqrt{y''} + 3x - 2y^2 = 3\sen(x)$
\end{enumerate}

\subsection*{Solucion Ejercicio 1}

\begin{enumerate}[a)]
\item Orden: 2, Grado: 2
\item Orden: 2, Grado: 1
\item Orden: 2, Grado: No definido (debido a la raíz cuadrada)
\end{enumerate}

\section*{Ejercicio 2}

Justifica cuáles de las siguientes ecuaciones diferenciales son lineales:

\begin{enumerate}[a)]
\item $y' + x^2 - 1 = y$
\item $3x\,y'' - 4y\,y' + y\sen(x) = 4x$
\item $y\,dx + (xy + x - 3y)\,dy = 0$
\item $y'' - 2y + 3 = e^x$
\item $4xy + \cos(y) = 5x$
\end{enumerate}

\subsection*{Solucion Ejercicio 2}
\begin{enumerate}[a)]
\item Lineal
\item No lineal (término $-4y\,y'$)
\item No lineal (término $xy\,dy$)
\item Lineal
\item No lineal (término $\cos(y)$)
\end{enumerate}

\section*{Ejercicio 3}

Comprueba que las funciones dadas son solución de sus correspondientes ecuaciones diferenciales:

\begin{enumerate}[a)]
\item $y = x + e^{-x}$ para la ecuación $y' + y = x + 1$
\item $y = 2e^{3x} - 5e^{4x}$ para la ecuación $y'' - 7y' + 12y = 0$
\item $y = e^x(2x + 1)$ para la ecuación $y'' = 2y' - y$
\item $x(t) = \sen(t) + \cos(t)$ para la ecuación $\dot{x}\cos(t) + x\sen(t) = 1$
\end{enumerate}

\subsection*{Solucion Ejercicio 3}
\begin{enumerate}[a)]
\item Derivando $y$ obtenemos $y' = 1 - e^{-x}$. Sustituyendo en la ecuación:
\[(1 - e^{-x}) + (x + e^{-x}) = x + 1.\] 
Por lo tanto, es solución.
\item Derivando $y$ obtenemos $y' = 6e^{3x} - 20e^{4x}$ y $y'' = 18e^{3x} - 80e^{4x}$. Sustituyendo en la ecuación:
\[(18e^{3x} - 80e^{4x}) - 7(6e^{3x} - 20e^{4x}) + 12(2e^{3x} - 5e^{4x}) = 0.\] 
Por lo tanto, es solución.
\item Derivando $y$ obtenemos $y' = e^x(2x + 3)$ y $y'' = e^x(2x + 5)$. Sustituyendo en la ecuación:
\[e^x(2x + 5) = 2e^x(2x + 3) - e^x(2x + 1).\] 
Por lo tanto, es solución.
\item Derivando $x(t)$ obtenemos $\dot{x}(t) = \cos(t) - \sen(t)$. Sustituyendo en la ecuación:
\[(\cos(t) - \sen(t))\cos(t) + (\sen(t) + \cos(t))\sen(t) = 1.\] 
Por lo tanto, es solución.
\end{enumerate}

\section*{Ejercicio 4}

Determina los valores de $m$ para los que la función
\[
f(x) = e^{mx}
\]
es solución de la ecuación
\[
y''' - 3y'' + 4y' - 2y = 0.
\]

\subsection*{Solucion Ejercicio 4}


Sea $f(x)=e^{mx}$.

\[
f'=me^{mx},\quad f''=m^2e^{mx},\quad f'''=m^3e^{mx}
\]

Sustituyendo:

\[
e^{mx}(m^3-3m^2+4m-2)=0
\]

\[
m^3-3m^2+4m-2=0
\]

\[
(m-1)(m^2-2m+2)=0
\]

\[
\boxed{m=1,\quad m=1\pm i}
\]
\section*{Ejercicio 5}

Verifica que la expresión paramétrica
\[
x(t) = \cos(t), \qquad y(t) = \sen(t)
\]
es solución de la ecuación
\[
x + y\,y' = 0.
\]

\subsection*{Solucion Ejercicio 5}
\[
x(t)=\cos t,\quad y(t)=\sen t
\]

\[
\frac{dy}{dx}=\frac{dy/dt}{dx/dt}=\frac{\cos t}{-\sen t}
\]

\[
x + yy' = \cos t + \sen t\left(\frac{\cos t}{-\sen t}\right)=0
\]
Por lo tanto, es solución.

\section*{Ejercicio 6}

Verifica que las siguientes expresiones implícitas son soluciones de las ecuaciones diferenciales asociadas:

\begin{enumerate}[a)]
\item $x^2 + y^2 - 4 = 0$ para la ecuación
\[
y' = -\frac{x}{y}.
\]

\item $\displaystyle \frac{y^2}{2} = \ln(1 + e^x) + C$ para la ecuación
\[
(1 + e^x)\, y\, y' = e^x.
\]
\end{enumerate}

\subsection*{Solucion Ejercicio 6}

\begin{enumerate}[a)]
\item Derivando implícitamente $x^2 + y^2 - 4 = 0$:
\[
2x + 2y y' = 0 \Rightarrow y' = -\frac{x}{y}
\]
Por lo tanto, es solución.

\item Derivando implícitamente $\displaystyle \frac{y^2}{2} = \ln(1 + e^x) + C$:
\[
y y' = \frac{e^x}{1 + e^x} \Rightarrow (1 + e^x) y y' = e^x
\]
Por lo tanto, es solución.
\end{enumerate}

\section*{Ejercicio 7}

Verifica que la función definida a trozos
\[
f(x) =
\begin{cases}
- x^2, & x < 0, \\
x^2, & x \ge 0,
\end{cases}
\]
es una solución de la ecuación diferencial
\[
x y' - 2y = 0.
\]

\subsection*{Solucion Ejercicio 7}

Para $\boxed{x < 0}$, $f(x) = -x^2$, por lo que $f'(x) = -2x$. Sustituyendo en la ecuación:
\[
x(-2x) - 2(-x^2) = -2x^2 + 2x^2 = 0.
\]
Para $\boxed{x \geq 0}$, $f(x) = x^2$, por lo que $f'(x) = 2x$. Sustituyendo en la ecuación:
\[
x(2x) - 2(x^2) = 2x^2 - 2x^2 = 0.
\]
Por lo tanto, es solución.

\section*{Ejercicio 8}

Dada la ecuación diferencial
\[
x^2 y'' - 4x y' + 6y = 0,
\]
completa los siguientes apartados:

\begin{enumerate}[a)]
\item Demuestra que tanto $y_1(x) = x^2$ como $y_2(x) = x^3$ son soluciones de la ecuación.
\item ¿Es $y(x) = y_1(x) + y_2(x)$ también solución de la ecuación?
\end{enumerate}

\subsection*{Solucion Ejercicio 8}
\begin{enumerate}[a)]
\item Para $y_1(x) = x^2$:
\[y_1' = 2x, \quad y_1'' = 2. \]
Sustituyendo en la ecuación:
\[x^2(2) - 4x(2x) + 6(x^2) = 2x^2 - 8x^2 + 6x^2 = 0.\]
Por lo tanto, $y_1(x)$ es solución.

Para $y_2(x) = x^3$:
\[y_2' = 3x^2, \quad y_2'' = 6x. \]
Sustituyendo en la ecuación:
\[x^2(6x) - 4x(3x^2) + 6(x^3) = 6x^3 - 12x^3 + 6x^3 = 0.\]
Por lo tanto, $y_2(x)$ es solución.
\item Sea $y(x) = y_1(x) + y_2(x) = x^2 + x^3$:
\[y' = 2x + 3x^2, \quad y'' = 2 + 6x. \]
Sustituyendo en la ecuación:
\[x^2(2 + 6x) - 4x(2x + 3x^2) + 6(x^2 + x^3) = 2x^2 + 6x^3 - 8x^2 - 12x^3 + 6x^2 + 6x^3 = 0.\]
Por lo tanto, $y(x)$ es también solución.
\end{enumerate}

\section*{Ejercicio 9}

Comprueba que la familia de funciones
\[
y = Cx + C^2
\]
es la solución general de la ecuación diferencial
\[
y = x y' + (y')^2.
\]
A continuación, determina el valor de $K$ para el que la función
\[
y = Kx^2
\]
sea una solución singular de la misma ecuación.

\subsection*{Solucion Ejercicio 9}
Derivando $y = Cx + C^2$:
\[y' = C. \]
Sustituyendo en la ecuación:
\[Cx + C^2 = x(C) + (C)^2.\]
Por lo tanto, es solución general.
Para $y = Kx^2$:
\[y' = 2Kx. \]
Sustituyendo en la ecuación:
\[Kx^2 = x(2Kx) + (2Kx)^2 \Rightarrow Kx^2 = 2Kx^2 + 4K^2x^2.\]
Igualando los coeficientes:
\[K = 2K + 4K^2 \Rightarrow 4K^2 + K = 0 \Rightarrow K(4K + 1) = 0.\]
Por lo tanto, $K = 0$ o $K = -\frac{1}{4}$. La solución singular es $K = -\frac{1}{4}$.

Nota: $K=0$ corresponde a una solución particular, que no es singular.

\section*{Ejercicio 10}

Sabiendo que la solución general de la ecuación diferencial
\[
x^2 y' \sen(y) = 1
\]
está dada por la relación
\[
\cos(y) = \frac{1}{x} + C,
\]
calcula la solución particular que tiende al valor $\displaystyle \frac{\pi}{2}$ cuando $x$ tiende a infinito.

\subsection*{Solucion Ejercicio 10}
Para que la solución particular tienda a $\frac{\pi}{2}$ cuando $x$ tiende a infinito, debemos tener:
\[\cos\left(\frac{\pi}{2}\right) = \frac{1}{\infty} + C \Rightarrow 0 = 0 + C \Rightarrow C = 0.\]
Por lo tanto, la solución particular es:
\[\cos(y) = \frac{1}{x} \Rightarrow y = \arccos\left(\frac{1}{x}\right).\]

\section*{Ejercicio 11}

Aplica el teorema de existencia y unicidad a la ecuación diferencial
\[
y' = y^{1/3},
\]
determinando los recintos del plano donde se puede asegurar que por cada punto pasa una única solución.

\subsection*{Solucion Ejercicio 11}

\[
\frac{\partial f}{\partial y}=\frac{1}{3}y^{-2/3}
\]

No es continua en $y=0$.

\[
\boxed{\text{El teorema de existencia y unicidad no se cumple en } y=0}
\]

\section*{Ejercicio 12}

Dado el problema de valor inicial
\[
\begin{cases}
y' - x\, y^{1/2} = 0 \\
y(0) = 0
\end{cases}
\]
comprueba que tanto $y = 0$ como $y = \dfrac{x^4}{16}$ son soluciones.  
¿Contradice este resultado el teorema de existencia y unicidad de soluciones en el origen de coordenadas?

\subsection*{Solucion Ejercicio 12}

Dos soluciones:

\[
y=0,\qquad y=\frac{x^4}{16}
\]

Ambas verifican la ecuación y pasan por $(0,0)$.

No contradice el teorema porque:

\[
\frac{\partial f}{\partial y}=\frac{x}{2\sqrt{y}}
\]

no es continua en $(0,0)$.

\section*{Ejercicio 13}

Determina la ecuación diferencial asociada a las siguientes soluciones generales:

\begin{enumerate}[a)]
\item $y - 4x - C = 0$
\item $y = C_1 e^x + C_2 e^{-x}$
\item $x^2 + C_1 y^2 = C_2 \quad ( \text{donde } C_1, C_2 > 0 )$
\end{enumerate}

\subsection*{Solucion Ejercicio 13}
\section*{Ejercicio 13}

Determina la ecuación diferencial asociada a las siguientes soluciones generales.

\begin{itemize}
\item[\textbf{a)}] $y - 4x - C = 0$

\[
y = 4x + C
\]

Derivando respecto de $x$:
\[
y' = 4
\]

Por tanto, la ecuación diferencial asociada es:
\[
\boxed{y' - 4 = 0}
\]

\item[\textbf{b)}] $y = C_1 e^x + C_2 e^{-x}$

Derivamos una primera vez:
\[
y' = C_1 e^x - C_2 e^{-x}
\]

Derivamos una segunda vez:
\[
y'' = C_1 e^x + C_2 e^{-x}
\]

Observamos que:
\[
y'' = y
\]

Luego, la ecuación diferencial asociada es:
\[
\boxed{y'' - y = 0}
\]

\item[\textbf{c)}] $x^2 + C_1 y^2 = C_2 \quad (C_1, C_2 > 0)$

Derivamos respecto de $x$:
\[
2x + 2C_1 y y' = 0
\]

Dividiendo entre 2:
\[
x + C_1 y y' = 0
\]

Despejamos $C_1$:
\[
C_1 = -\frac{x}{y y'}
\tag{1}
\]

Derivamos nuevamente la ecuación $x + C_1 y y' = 0$:
\[
1 + C_1 \frac{d}{dx}(y y') = 0
\]

Derivando el producto:
\[
\frac{d}{dx}(y y') = (y')^2 + y y''
\]

Sustituyendo:
\[
1 + C_1\left((y')^2 + y y''\right) = 0
\]

Despejamos $C_1$:
\[
C_1 = -\frac{1}{(y')^2 + y y''}
\tag{2}
\]

Igualando $C_1$ de ambas ecuaciones:
\[
-\frac{x}{y y'} = -\frac{1}{(y')^2 + y y''}
\]

Simplificando:
\[
\frac{x}{y y'} = \frac{1}{(y')^2 + y y''}
\]

\[
x\left((y')^2 + y y''\right) = y y'
\]

\[
x (y')^2 + x y y'' = y y'
\]

\[
\boxed{x y y'' + x (y')^2 - y y' = 0}
\]

\end{itemize}

\section*{Ejercicio 14}

Escribe la ecuación diferencial asociada a una solución general representada por una familia de curvas para la que se verifica que, en todo punto, la pendiente de la recta tangente a cada curva es igual a la distancia de dicho punto al origen de coordenadas.

\subsection*{Solucion Ejercicio 14}

Sea $y = f(x)$ una curva de la familia. La pendiente de la recta tangente en un punto $(x, y)$ es $f'(x)$. La distancia del punto al origen es $\sqrt{x^2 + y^2}$. Por lo tanto, la ecuación diferencial asociada es:
\[y' = \sqrt{x^2 + y^2}.\]

\end{document}